\documentclass[11pt,a4paper]{article}
\usepackage[T1]{fontenc}
\usepackage[utf8]{inputenc}
\usepackage{amsmath,amssymb,amsthm}
\usepackage[margin=1in]{geometry}
\usepackage[colorlinks=true,linkcolor=blue,urlcolor=blue]{hyperref}
\theoremstyle{plain}
\newtheorem{theorem}{Theorem}
\newtheorem{lemma}[theorem]{Lemma}
\newtheorem{proposition}[theorem]{Proposition}
\newtheorem{corollary}[theorem]{Corollary}
\theoremstyle{definition}
\newtheorem{remark}[theorem]{Remark}

\title{Recovering the unwritten recurrences of the table OEIS A237850}
\author{Adrian Perez Fontelles\\ \small Independent researcher}
\date{2 September 2026}
\begin{document}
\maketitle

\begin{abstract}
OEIS A237850 is a two-dimensional table $T(n,k)$ whose empirical recurrences are, for several of
its lines, given only as an ORDER --- ``k=2: [order 29],'' and the like --- with the
coefficients in a linked file rather than in the entry. Those conjectures can still be settled,
and the recurrences recovered. A line of the table is a fixed-width or fixed-height array
count, hence a walk count on finitely many vertices, hence satisfies some monic recurrence of
order at most that number; Berlekamp--Massey on twice as many exact terms returns the MINIMAL
one. Where its order equals the order the entry states --- and where the same holds after the
first few terms are dropped --- any recurrence of that order the line satisfies has a
characteristic polynomial that is a multiple of the minimal one and of equal degree, hence is
that polynomial. This note settles column 2, column 3.
\end{abstract}

\noindent\small 2020 Mathematics Subject Classification. 05A15, 11B37, 15A18.\normalsize

\section{The table and the unwritten conjectures}

OEIS A237850 is ``T(n,k) = Number of (n+1) X (k+1) 0..2 arrays with the maximum minus the lower median of every 2 X 2 subblock equal.''. It is read by antidiagonals and begins
\[
81, 369, 369, 1697, 3419, 1697, 8451, 31661,\ \dots
\]
The entry carries, among its empirical claims:
\begin{quote}\small
k=2: [order 29],\\
k=3: [order 70].
\end{quote}
As of the ``Last modified'' line on the live entry (Apr 13 2022 20:18:08, revision 6) these are still
recorded as empirical, and the recurrences themselves are not in the entry text.

\section{Why a line of the table is a walk count}

Fix the index that the line fixes. The entry defines $T(n,k)$ as a number of arrays of a shape
determined by $n$ and $k$, subject to a condition local to a bounded window of consecutive
lines. Substituting the fixed index into the entry's own wording turns the definition into an
ordinary one-parameter array count, and for such a count the admissible windows are the
vertices of a finite digraph and an array is a walk. So the line is a walk count on $S$
vertices, and by the Cayley--Hamilton theorem it satisfies a monic integer recurrence of order
at most $S$.

\section{Recovering the recurrences}

\begin{lemma}\label{lem:bm}
A sequence known to satisfy some linear recurrence of order at most $S$ is determined, with its
minimal recurrence, by its first $2S$ terms; Berlekamp--Massey applied to them returns that
minimal recurrence exactly.
\end{lemma}

\subsection*{Column $k=2$}
Substituting the fixed index into the entry's wording gives an ordinary one-parameter array
count, a walk on $S=66$ vertices. Berlekamp--Massey on $2S=132$ exact terms
returns a minimal recurrence of order $29$ --- the order the entry states --- with
integer coefficients
\[
(c_1,\dots,c_{29})=(19,\ -26,\ -1354,\ 6471,\ 28813,\ -231046,\ -109423,\ 3526315,\ -3576784,\ -26974056,\ 52871335,\ 103996087,\ -324210864,\ -147836704,\ 1068205642,\ -281417288,\ -1965277664,\ 1461313160,\ 1869996168,\ -2374039488,\ -570178496,\ 1797323328,\ -380875136,\ -564139008,\ 295181312,\ 22262784,\ -45019136,\ 9666560,\ -491520).
\]
so that $a(m)=\sum_{i=1}^{29}c_i\,a(m-i)$ for every $m$. Removing the first
$29+4$ terms and repeating gives order $29$ again, which is what the
identification below needs.

\subsection*{Column $k=3$}
Substituting the fixed index into the entry's wording gives an ordinary one-parameter array
count, a walk on $S=171$ vertices. Berlekamp--Massey on $2S=342$ exact terms
returns a minimal recurrence of order $70$ --- the order the entry states --- with
integer coefficients
\[
(c_1,\dots,c_{70})=(29,\ 212,\ -11594,\ 16048,\ 1851803,\ -8142699,\ -161957175,\ 1018196013,\ 8759605120,\ -70791362382,\ -308024111873,\ 3226028798774,\ 6991536735908,\ -103773108838562,\ -88876774693040,\ 2459509894134199,\ -29718656767235,\ -44162259405169850,\ 27903095438170592,\ 612375987544108578,\ -675647143703665928,\ -6647034900341286560,\ 9872761858009910305,\ 57019345037769166467,\ -103561560629448744210,\ -389046410915450979941,\ 826802405921981697834,\ 2119446784165876553094,\ -5167503124808330098395,\ -9230550030878199073073,\ 25695519904846129976904,\ 32082555344956204608426,\ -102696948207908747745380,\ -88470689861335223816157,\ 332106805798664891973145,\ 191092255163057612377480,\ -872721691626771756306528,\ -314621596281858441727192,\ 1868130830625283825144686,\ 369683330602431932092578,\ -3260154604117915151179776,\ -244740435375417184187156,\ 4635396789184137041538152,\ -75988235571601561857080,\ -5358219727367962615339600,\ 441055859959968374913904,\ 5016597735990622931562176,\ -635886941330459459715136,\ -3782944029707679206804928,\ 577437197600364143509632,\ 2279773917855028474530560,\ -369836611493021336015872,\ -1086273458889201801969664,\ 171111626468250107523072,\ 403231504807605571710976,\ -56746331075032955756544,\ -114221216369879807819776,\ 13092510454272728465408,\ 23968849518126804566016,\ -1987179770742385934336,\ -3567823762819877175296,\ 179962855775486869504,\ 352913709341046996992,\ -8252694928281305088,\ -20952256695301570560,\ 191144368948641792,\ 629102097350000640,\ -9645259864670208,\ -6772884252917760,\ 238654152769536).
\]
so that $a(m)=\sum_{i=1}^{70}c_i\,a(m-i)$ for every $m$. Removing the first
$70+4$ terms and repeating gives order $70$ again, which is what the
identification below needs.

\begin{theorem}
For each line above, the displayed recurrence holds for every index, and it is the recurrence
the entry records.
\end{theorem}

\begin{proof}
It holds by Lemma~\ref{lem:bm}. For the identification, the entry asserts that the line
satisfies SOME recurrence of the stated order, possibly only from some index on. Let $q$ be its
characteristic polynomial and $q_{\min}$ the minimal polynomial of the tail. Then
$q_{\min}\mid q$, and the second Berlekamp--Massey run --- on the line with its first few
terms removed --- shows $\deg q_{\min}$ equals the stated order, which is $\deg q$. Both are
monic, so $q=q_{\min}$.
\end{proof}

\section{Verification}

Three checks.

First, each line's model was compared against the entry's own published values for that line,
taken from the antidiagonal DATA read in either orientation and from the ``Table starts''
block, and agrees at every available term. This is the only point at which reading the entry's
English enters, and a misreading gives different counts.

Second, each recovered recurrence was evaluated directly on the model's values, with no
Berlekamp--Massey involved, and holds at every index tested.

Third, each order was computed twice by independent routes --- modulo a $61$-bit prime, where
every intermediate is a machine word, and again in exact rational arithmetic --- and once more
on the line with its first few terms dropped, which is what the identification requires. All
agree.

\begin{thebibliography}{9}
\bibitem{oeis} The OEIS Foundation, \emph{The On-Line Encyclopedia of Integer
Sequences}, \url{https://oeis.org}, 2026. Sequence A237850.
\bibitem{massey} J.~L.~Massey, Shift-register synthesis and BCH decoding, \emph{IEEE
Trans. Inform. Theory} \textbf{15} (1969), 122--127.
\bibitem{stanley} R.~P.~Stanley, \emph{Enumerative Combinatorics}, Volume 1, 2nd ed.,
Cambridge University Press, 2011. (Transfer-matrix method, Section 4.7.)
\bibitem{horn} R.~A.~Horn and C.~R.~Johnson, \emph{Matrix Analysis}, 2nd ed., Cambridge
University Press, 2013.
\end{thebibliography}

\end{document}
